\section{¿Qué pasaría si...? Ideas especulativas para el proyecto}

\subsection{¿Qué pasaría si Apophis tuviera masa comparable a la de la Luna?}

Simular el encuentro con una masa artificial mucho mayor para Apophis, del orden de la masa lunar. Estudiar a partir de qué escala de masa la gravedad del asteroide empieza a perturbar apreciablemente la órbita de la Tierra durante el encuentro. Evaluar además si existe una masa crítica para la cual Apophis podría quedar capturado temporal o permanentemente por el sistema Tierra--Sol.

\subsection{¿Qué pasaría si el encuentro ocurriera exactamente en el plano de la eclíptica?}

Forzar $I = 0^{\circ}$ en los elementos orbitales de Apophis y repetir la simulación. Analizar si la distancia mínima de aproximación cambia de manera significativa y explorar si existe alguna inclinación crítica que maximice o minimice la distancia de encuentro.

\subsection{¿Qué pasaría si Júpiter no existiera?}

Re-simular los últimos 100 años de la órbita de Apophis eliminando a Júpiter del sistema. Comparar el estado orbital resultante con el escenario nominal y determinar si el encuentro de 2029 habría ocurrido con una geometría distinta. Esta prueba permite cuantificar cuánto influye Júpiter como perturbador dominante del sistema solar interior en este caso particular.

\subsection{¿Qué pasaría si el encuentro hubiera ocurrido del lado nocturno de la Tierra?}

Tener en cuenta la rotación terrestre durante el paso cercano y calcular qué punto geográfico de la Tierra quedaría bajo Apophis en el instante de máxima aproximación. A partir de ello, identificar desde qué regiones o ciudades del mundo el asteroide sería visible a simple vista y comparar ese escenario con la geometría real del encuentro.

\subsection{¿Qué pasaría si Apophis pasara exactamente a la distancia de los satélites geoestacionarios?}

Ajustar artificialmente la órbita de Apophis para que la distancia mínima a la Tierra sea exactamente $42{,}164\,\mathrm{km}$. Calcular la deflexión angular asociada a ese paso cercano y el cambio resultante en su período orbital heliocéntrico. Como extensión, evaluar si después de esa perturbación seguiría clasificándose como un asteroide tipo Aten.